% 原书物理页=印刷页12--14；2.1 全节（末尾类公式显示式位于第 14 页顶部）.
\chapter{作用、半直积与经典群}\label{chap:fg-2}
\section{群在集合上的作用}\label{sec:fg-2-1}
\begin{definition}\label{def:fg-2-1-1}
设 $G$ 是一个群，$X$ 是一个集合. $G$ 在 $X$ 上的一个\newterm{fg-group-action}{（左）作用}{(left) action}是一个映射
\[
\begin{array}{rcl}
G\times X &\longrightarrow& X\\
(g,x) &\longmapsto& g.x
\end{array}
\]
它满足下列性质：
\begin{itemize}
\item 对任意 $x\in X$，有 $e.x=x$.
\item 对任意 $x\in X$ 与任意 $(g,h)\in G\times G$，有 $g.(h.x)=(gh).x$.
\end{itemize}
\end{definition}
这里同样可以用几种方式来记这个作用：$g\cdot x$、$g*x$，或简单地写作 $gx$，有时也写作 ${}^{g}x$. 给出 $G$ 在 $X$ 上的一个作用，等价于给出一个从 $G$ 到 $X$ 的置换群 $\mathrm{S}_X$ 的同态 $\tau$：只需规定 $\tau(g)(x)=g.x$，并注意 $\tau(g^{-1})=(\tau(g))^{-1}$. 也称 $G$（通过 $\tau$）（左）作用在 $X$ 上. 同样可以定义 $G$ 在 $X$ 上的右作用：$(g,x)\mapsto xg$，此时要求 $x(gh)=(xg)h$.
\begin{numpara}[轨道]\label{np:fg-2-1-2}
设 $G$ 是一个作用在集合 $X$ 上的群（当不特别说明时，意指``左作用''）. 若 $A$（相应地，$Y$）是 $G$（相应地，$X$）的一个子集，则把 $X$ 中形如 $ay$（$a\in A$，$y\in Y$）的元素所成的集合记作 $AY$：
\[
AY=\{x\in X,\ \exists(a,y)\in A\times Y,\ x=a.y\}.
\]
对任意 $x\in X$，$X$ 的子集 $Gx:=G\{x\}=\{gx,\ g\in G\}$ 称为 $x$ 在 $G$ 下的\newterm{fg-orbit}{轨道}{orbit}（或 $G$-轨道）. 此时，$Gx=Gy$ 当且仅当 $y\in Gx$，并且各 $G$-轨道所成的集合 $G\setminus X$ 构成 $X$ 的一个划分. 若 $G$ 在 $X$ 上的作用是右作用，则把它记作 $X/G$.
\end{numpara}
称 $G$ \newterm{fg-transitive-action}{传递地}{transitively}作用在一个集合 $X$ 上，如果 $G\setminus X$ 只有一个元素. 例如，$G$ 在它的每条轨道上都传递地作用.
\begin{numpara}[不动点与稳定子]\label{np:fg-2-1-3}
设 $x$ 是 $X$ 的一个点，$H$ 是 $G$ 的一个子群. 若 $\forall h\in H,\ hx=x$，则称 $x$ 被 $H$ 固定，并把 $X$ 中被 $H$ 固定的点所成的集合记作 $X^H$. 反过来，对任意 $x\in X$，把子群 $G_x:=\{g\in G,\ gx=x\}$ 称为 $x$ 的\newterm{fg-stabilizer}{稳定子}{stabilizer}（或\newterm{fg-fixator}{固定子}{fixator}）. 称 $G$ \newterm{fg-faithful-action}{忠实地}{faithfully}作用在 $X$ 上，如果 $\bigcap_{x\in X}G_x=\{1\}$，也就是说同态 $\tau$ 是单射. 更一般地，对 $X$ 的任意子集 $Y$，把 $G$ 的子群 $\bigcap_{x\in Y}G_x$ 称为 $Y$ 的固定子 $\mathit{Fix}_G(Y)$，它由那些逐点固定 $Y$ 的元素组成. 一般地说，$G$ 中整体保持 $Y$ 稳定的元素所成的集合 $\{g\in G,\ gY\subset Y\}$ 只是 $G$ 的一个子幺半群（即对群的运算稳定且包含 $1$ 的子集）. 不过，它常常（例如当 $Y$ 有限时）与 $G$ 的子群 $G_Y:=\{g\in G,\ gY=Y\}$ 重合，此时就把后者称为 $Y$ 的稳定子.
\end{numpara}
\begin{numpara}[例]\label{np:fg-2-1-4}
\begin{enumerate}[label=\roman*)]
\item 设 $H$ 是 $G$ 的一个子群. 此时 $H$ 通过平移左作用在 $G$ 上：$(h,g)\mapsto hg$，而一条轨道 $Hg$ 恰好是模 $H$ 的一个右陪集. 因此，（\ref{np:fg-1-1-5} 段与 \ref{np:fg-2-1-2} 段中的）两种 $H\setminus G$ 记号是相容的. 对于 $H$ 在 $G$ 上的右平移作用，对 $G/H$ 亦有类似结论. 对这两个作用中的每一个，任何元素的固定子都仅是 $\{1\}$.
\item 由同态 $\tau=\operatorname{Int}$ 给出的 $G$ 在 $X=G$ 上的（左）作用 $(g,x)\mapsto gxg^{-1}$ 称为 $G$ 在 $G$ 上的共轭作用. 若 $G$ 的两个元素位于同一条轨道中，则称它们是共轭的. 轨道 $C_x:=\{gxg^{-1},\ g\in G\}$ 称为 $G$ 的（元素的）\newterm{fg-conjugacy-class}{共轭类}{conjugacy class}. $G$ 的共轭类所成的集合记作 $C\ell(G)$，或记作 $G/\mathit{ad}\,G$；它的基数 $c\ell(G)$ 称为 $G$ 的类数. 对 $G$ 的任意子集 $Y$，$Y$ 的固定子 $Z_G(Y)=\{g\in G,\ \forall y\in Y,\ gy=yg\}$ 称为 $Y$（在 $G$ 中）的\newterm{fg-centralizer}{中心化子}{centralizer}（或交换子）. 子群 $N_G(Y)=\{g\in G,\ gY=Yg\}$ 称为 $Y$（在 $G$ 中）的\newterm{fg-normalizer}{正规化子}{normalizer}. $G$ 自身的中心化子就是 $G$ 的中心 $Z(G)=Z_G(G)$.
\item 设 $S(G)$ 是 $G$ 的全体子群所成的集合. 对任意 $g\in G$ 与 $H\in S(G)$，$H'=gHg^{-1}$ 仍是 $G$ 的子群，这就赋予 $S(G)$ 一个 $G$ 的作用. 若 $G$ 的两个子群位于同一条轨道中，则称它们是共轭的. 这些轨道称为 $G$ 的子群的共轭类. 对 $H\in S(G)$，$H$ 的固定子不是别的，正是它的正规化子 $N_H(G)$\footnote{译注：原文如此，按定义应为 $N_G(H)$.}；它是 $G$ 的使 $H$ 正规的最大子群. 例如，$H\triangleleft G\iff N_H(G)=G$.
\end{enumerate}
\end{numpara}
\begin{numpara}\label{np:fg-2-1-5}
设 $G$ 是一个作用在集合 $X$ 上的群. 对任意 $x\in X$，映射 $g\mapsto gx$ 诱导出一个从 $G/G_x$ 到 $G.x$ 的双射（其中 $1$ 所在的陪集被送到 $x$）. 特别地，
\[
|G.x|=[G:G_x].
\]
若把 $x$ 换成 $x'=gx\in G.x$，则 $G_x$ 就被换成子群 $G_{x'}=gG_xg^{-1}$，反之亦然. 因此，给出一个 $G$ 传递作用于其上的 $X$，等价于给出 $G$ 中某个稳定子的共轭类. 例如，对任意域 $K$，由仿射直线 $X=\mathbb{A}^1(K)$ 的相似变换 $x\mapsto ax+b$（$a\in K^\times$，$b\in K$）组成的群 $G=\mathit{Aff}^1(K)$ 传递地作用在 $X$ 上；对任意 $x\in X$，$x$ 的稳定子是以 $x$ 为中心的位似所成的子群，从而可以把 $X$ 等同于子群 $\mathrm{GL}_1(K)=K^\times$ 在 $G$ 中的共轭类.
\end{numpara}
\begin{numpara}\label{np:fg-2-1-6}
设 $G$ 是一个作用在有限集合 $X$ 上的有限群，$x_1,\ldots,x_m$ 是 $X$ 在 $G$ 作用下各轨道的一个代表元系. 换言之，$m$ 是 $G\setminus X$ 的基数，并且 $X=\bigcup_{i=1,\ldots,m}G.x_i$（不交并）. 于是 $|X|=\sum_{i=1,\ldots,m}[G:G_{x_i}]=\sum_{i=1,\ldots,m}\frac{|G|}{|G_{x_i}|}$，从而
\[
\frac{|X|}{|G|}=\sum_{i=1}^{m}\frac{1}{|G_{x_i}|}.
\]
\end{numpara}
我们把这个``类公式''应用于 $G$ 通过共轭作用在其自身上的作用. 此时，$m=c\ell(G)$ 是 $G$ 的共轭类数，每个 $G_{x_i}$ 都是 $x_i$ 在 $G$ 中的中心化子 $Z_G(x_i)$（例如，$x_1=1$ 所在的类仅由一个元素组成，且满足 $Z_G(1)=G$），于是得到
\[
1=\sum_{i=1}^{c\ell(G)}\frac{1}{|Z_G(x_i)|}.
\]
